\section{Separated prime supports and discounted certificates}
\label{ss:section}

This section proves a uniform coefficient-two asymptotic bound for sources
whose prime factors lie in two separated ranges. The large-prime range may
start at $M^\alpha$ for any fixed $\alpha>0$. A quantitative extension also
allows a specified exponent tending to zero. The intermediate prime range
remains excluded.

All logarithms are natural unless indicated otherwise. For a prime set $Q$,
write $Q^\Pi=\prod_{q\in Q}q$ and
$\vartheta(Q)=\prod_{q\in Q}(1-1/q)$, with empty products equal to one.
An integer is $Q$-smooth when all its prime factors belong to $Q$;
one is included. Intervals at real lengths are half-open intervals
$(x,x+T]$ with $x\ge0$, with their integer points understood.
The ambient integer $M$ is retained when sources are deleted.

\subsection{Precise analytic inputs}

We recall the following previously proved statements, including the
coefficient in the diffuse comparison. Their proofs, with partial
prime-power carriers and uniform interval errors, are in \cite{Synthesis}.
The new arguments below use the coefficients as stated.

\begin{lemma}[Uniform sieve and diffuse comparison inputs]
\label{ss:inputs}
There is an absolute $C>0$ such that
\[
 e(u)=C\bigl((u/2)^{-u/4}+2^{1-u/2}\bigr),\qquad u\ge4,
\]
is decreasing to zero and has the following property. If $Q$ is any
subset of the primes at most $z\ge2$ and an interval $J$ has real length
$T\ge z^u$, then
\[
 \left|\#\{m\in J:(m,Q^\Pi)=1\}-T\vartheta(Q)\right|
 \le e(u)T\vartheta(Q).
 \tag{S}
\]
For weights $0\le w_p\le\eta$, put
$V(m)=\sum_pw_pv_p(m)$ and $\mu=\sum_pw_p/(p-1)$.
Let $\ell>\eta$, $r=\lfloor\ell/\eta\rfloor\ge8$,
$j=\lfloor r/2\rfloor$, and $\kappa=r-j$.
Whenever $e(r)<1$, every level $s\ge0$ and every pair of a prefix and
an interval of the same real length $T$ satisfy
\begin{equation}
 \#\{a\le T:V(a)>s+\ell\}
 \le
 \frac{(1+e(\kappa))\mu+\eta2^{1-j}}
      {(\ell-\eta)(1-e(r))}
 \#\{b\in J:V(b)>s\}.
 \label{ss:diffuse-ratio}
\end{equation}
Finally, if the support is contained in $p\le M$, then
\begin{equation}
 \sum_{a=1}^M V(a)\ge M(\mu-3\eta).
 \label{ss:small-mean-error}
\end{equation}
\end{lemma}

For clarity, the carrier proof of \eqref{ss:diffuse-ratio} bounds the
source count by a common carrier sum times
$((1+e(\kappa))\mu+\eta2^{1-j})/(\ell-\eta)$, and bounds the target
count below by that sum times $1-e(r)$. Thus retaining their quotient
gives exactly the displayed statement.
We also use the classical uniform bounds
\begin{equation}
 \pi(T)\le C_0T/\log T,\qquad
 \vartheta(Q)\ge c_0/\log Y
 \quad(Q\subseteq\{p:p\le Y\},\ Y\ge2),
 \label{ss:prime-inputs}
\end{equation}
with absolute positive constants. These follow from Chebyshev and
Mertens estimates \cite{GranvilleSound,TaoMertens}.

\subsection{A strict budget for arbitrary pure-core sources}

\begin{proposition}[Uniform pure-core contraction]
\label{ss:pure-core}
Let $Y=Y(M)\ge2$ be prescribed with $\log Y=o(\sqrt{\log M})$.
For every $\varepsilon>0$ there is $\rho_\varepsilon<1$ such that,
for all sufficiently large $M$, the following holds uniformly in
$Q\subseteq\{p:p\le Y\}$, weights $0\le w_q\le1$, distinct
$Q$-smooth sources $A_Q\subseteq[2,M]$, and positive $M$-term intervals:
\begin{equation}
 \sum_{a\in A_Q}(V(a)-2-\varepsilon)_+
 \le \rho_\varepsilon
       \sum_{b\in I}(V(b)-1)_+,
 \qquad V(m)=\sum_{q\in Q}w_qv_q(m).
 \label{ss:pure-contraction}
\end{equation}
No bound on the number of prime factors of a pure-core source is assumed.
\end{proposition}

\begin{proof}
It suffices to treat $0<\varepsilon\le1$.
Set $\ell=1+\varepsilon$ and
\[
 \rho_\varepsilon=\frac{1+\varepsilon/2}{1+\varepsilon}<1.
\]
Choose $\eta>0$ so small that $\eta\le\varepsilon/24$, $r\ge8$,
$e(r)<1$, and
\begin{equation}
 \frac{(1+e(\kappa))(1+\varepsilon/4)+\eta2^{1-j}}
      {(\ell-\eta)(1-e(r))}
 \le\rho_\varepsilon,
 \label{ss:strict-choice}
\end{equation}
where $r,j,\kappa$ are as in Lemma~\ref{ss:inputs}.
Such a choice exists: as $\eta\downarrow0$, the left side tends to
$(1+\varepsilon/4)/(1+\varepsilon)$, strictly below $\rho_\varepsilon$.
All these parameters depend only on $\varepsilon$.

We first record a uniform smooth-number bound. Let $\Psi(M,Y)$ count
positive integers at most $M$ all of whose prime factors are at most $Y$.
Write $N=\log M$, $y=\log Y$, and $\sigma=1/(4y)$.
For $p\le Y$, put $z_p=p^{-1+\sigma}$. Then
$z_p\le e^{1/4}/p$ and $z_p/(1-z_p)\le4/p$. Consequently
\[
 \prod_{p\le Y}(1-p^{-1+\sigma})^{-1}
 \le\exp\left(4\sum_{p\le Y}\frac1p\right)\le Y^4.
\]
Rankin's inequality therefore gives
\begin{equation}
 \Psi(M,Y)\le
 M\exp\left(-\frac{N}{4y}+4y\right).
 \label{ss:smooth-count}
\end{equation}
For every fixed $B\ge0$, it follows that
\begin{equation}
 \frac{\Psi(M,Y)\log_2M}{M/Y^B}
 \le(\log_2M)
       \exp\left(-\frac N{4y}+(B+4)y\right)\longrightarrow0.
 \label{ss:smooth-small}
\end{equation}
Indeed $y=o(\sqrt N)$ makes $N/y$ dominate both $y$ and $\log N$.
The bounds are independent of the prime subset and the weights.

Put $\mu=\sum_{q\in Q}w_q/(q-1)$ and
$H(V)=\sum_{b\in I}(V(b)-1)_+$.
There are three cases.

\emph{All weights are at most $\eta$, and $\mu\le1+\varepsilon/4$.}
Equations \eqref{ss:diffuse-ratio} and \eqref{ss:strict-choice} give
a level-count ratio at most $\rho_\varepsilon$.
Integrating over $s\ge1$ yields
\[
 \sum_{a\le M}(V(a)-2-\varepsilon)_+
 \le\rho_\varepsilon H(V).
\]
Distinctness permits restriction to $A_Q$.
This case includes the empty support and zero weights.

\emph{All weights are at most $\eta$, and $\mu\ge1+\varepsilon/4$.}
For large $M$, all $q\le Y$ are below $M$.
Every interval prime-power count is at least its prefix count, so
\[
 \begin{split}
 H(V)&\ge\sum_{b\in I}V(b)-M
 \ge M(\mu-3\eta-1)
 \ge\varepsilon M/8.
 \end{split}
\]
For $a\le M$, the source load satisfies
$V(a)\le\Omega(a)\le\log_2M$. The source hinge is at most
$\Psi(M,Y)\log_2M=o(M)$ by
\eqref{ss:smooth-small} with $B=0$. Its ratio to $H(V)$ is eventually
at most $\rho_\varepsilon$.

\emph{Some weight $w_q$ is at least $\eta$.}
Put $j_q=\lceil2/w_q\rceil$ and $B=\lceil2/\eta\rceil$.
Every multiple of $q^{j_q}$ has $V\ge2$, while
$q^{j_q}\le Y^B$ and $M/Y^B\to\infty$. Hence
\[
 H(V)\ge\left\lfloor\frac M{q^{j_q}}\right\rfloor
 \ge\frac M{2Y^B}
\]
for large $M$. The source hinge is again at most
$\Psi(M,Y)\log_2M$, which is $o(M/Y^B)$ by
\eqref{ss:smooth-small}. Its ratio to $H(V)$ is eventually at most
$\rho_\varepsilon$.
The thresholds in all three cases are uniform, proving
\eqref{ss:pure-contraction}.
\end{proof}

\subsection{A pointwise certificate with controlled multiplicity}

Let $Q,R_0$ be disjoint finite prime sets, and let $\mathcal D$ be a finite
set of integers $d\le M$ with the following properties:
\[
 d\text{ is }(Q\cup R_0)\text{-smooth},\qquad
 \Omega_{R_0}(d)\le t,\qquad
 r\mid d,\ r\in R_0\ \Longrightarrow\ r\le T_d:=M/d.
 \tag{D}
\]
Here $t\ge0$ is an integer, and $\Omega_{R_0}$ counts occurrences with
multiplicity. The number of $Q$-prime occurrences is unrestricted.
For an integer $L\ge1$, set $K_d(b)=0$ when $d\nmid b$, and otherwise set
\begin{equation}
 K_d(b)=\mathbf1_{(b/d,Q^\Pi)=1}
       \left(1-\frac1L
          \sum_{\substack{r\in R_0\\r\le T_d}}\mathbf1_{r\mid b/d}\right).
 \label{ss:discounted}
\end{equation}
Define
\begin{align}
 C_1(t,L)&=\binom{2t+L-1}{t},\notag\\
 C_2(t,L)&=(t+L)\sum_{j=0}^{L-1}
                  \binom{t+L-1}{j}(t+1)^j,\notag\\
 C_*(t,L)&=\min\{C_1(t,L),C_2(t,L)\}.
 \label{ss:multiplicity}
\end{align}

\begin{lemma}[Discounted pointwise capacity]
\label{ss:capacity}
Let $X(m)=\sum_pw_pv_p(m)$ have nonnegative weights on any finite prime
support containing $Q\cup R_0$, and put $h_d=(X(d)-1)_+$.
Then, for every positive integer $b$,
\begin{equation}
 \sum_{d\in\mathcal D}h_dK_d(b)
 \le C_*(t,L)(X(b)-1)_+.
 \label{ss:pointwise}
\end{equation}
\end{lemma}

\begin{proof}
Suppose a positive $K_d(b)$ exists, and choose the numerically least
positive carrier $d_0$. Then $T_0=M/d_0$ is the largest threshold
among positive carriers. Put
\[
 S=\{r\in R_0:r\le T_0,\ r\mid b\},\qquad k=|S|.
\]
Every member of $S$ divides $d_0$ or $b/d_0$.
The first class has at most $t$ distinct members.
The second has at most $L-1$ members, since positivity in
\eqref{ss:discounted} means that its integer sum is less than $L$.
Consequently $k\le t+L-1$.

Every positive carrier has the same $Q$-part, namely the full $Q$-smooth
part of $b$. By (D), its $R_0$-part is a monomial on $S$ of total
degree at most $t$. There are $\binom{k+t}{t}\le C_1(t,L)$ possible
monomials, including the constant monomial. For $k=0$ this count is one.
The common $Q$-part and the monomial determine the carrier.

For the sharper count, order $S$ increasingly.
For a positive carrier $d$, the primes in $S$ at most $T_d$ form a
prefix, with at most $k+1$ choices. Outside that prefix, their exponents
in $d$ are zero by (D). Inside it, all but at most $L-1$ primes must
occur to their full exponent $v_r(b)$: every deficiency is counted
by $r\mid b/d$ in \eqref{ss:discounted}.
Choose the exceptional primes and their exponents, each in
$\{0,\ldots,t\}$. All other exponents are fixed by $b$. Thus the number
of positive carriers is at most
\[
 (k+1)\sum_{j=0}^{L-1}\binom{k}{j}(t+1)^j\le C_2(t,L).
 \]
Choices violating another restriction only enlarge this upper bound.

For each positive term, $K_d(b)\le1$ and $d\mid b$ implies
$h_d\le(X(b)-1)_+$. All other weighted terms are nonpositive.
Either multiplicity bound therefore proves \eqref{ss:pointwise};
if no term is positive the inequality is immediate.
\end{proof}

\subsection{Uniform interval averages of the certificate}

\begin{lemma}[A positive average despite a large reciprocal sum]
\label{ss:average}
Suppose $Q\subseteq\{p:p\le Y\}$, all occupied carriers have $T_d\ge Z$,
and $Z\to\infty$. Allow $t,L,Y,Z$ to depend on $M$.
Assume uniformly over the admitted data that
\begin{equation}
 \frac{\mu_0}{L}\le\frac14,\qquad
 \mu_0=\sum_{r\in R_0}\frac1r,\qquad
 \log Y=o(\log Z),\qquad
 \frac{(\log Y)^2}{L\log Z}\longrightarrow0.
 \label{ss:average-conditions}
\end{equation}
Then, for all sufficiently large $M$, uniformly in the interval shift
and $d\in\mathcal D$,
\begin{equation}
 G_d:=\sum_{b\in I}K_d(b)
 \ge\frac12\vartheta(Q)T_d.
 \label{ss:positive-average}
\end{equation}
\end{lemma}

\begin{proof}
Let $C_Q(J)$ count integers in a real interval $J$ coprime to $Q^\Pi$.
For $T=M/d$ and $J=(x/d,(x+M)/d]$, exact division gives
\begin{equation}
 G_d=C_Q(J)-\frac1L
      \sum_{\substack{r\in R_0\\r\le T}}C_Q(J/r).
 \label{ss:average-identity}
\end{equation}
Division by $r\in R_0$ preserves the $Q$-coprimality condition.

Fix an absolute $\kappa\ge4$ with $e(\kappa)\le1/16$ and put
$U=T/Y^\kappa$. By \eqref{ss:average-conditions}, for large $M$,
$T\ge Y^\kappa$ and $U\ge\sqrt T$, uniformly for $T\ge Z$.
Apply (S) to the base count and every subtracted term with $r\le U$.
Their combined contribution is at least
\[
 \begin{split}
 \vartheta(Q)T\left[
  1-\frac{\mu_0}{L}
       -e(\kappa)\left(1+\frac{\mu_0}{L}\right)\right]
 \ge\frac{43}{64}\vartheta(Q)T.
 \end{split}
\]
For the remaining terms, discard coprimality and use
$C_Q(J/r)\le T/r+1$. Chebyshev's bound and partial summation give
\[
 \begin{split}
 \sum_{U<r\le T}\frac1r
 &\le\frac C{\log T}
       +C\int_U^T\frac{du}{u\log u}\\
 &\le C'\frac{1+\kappa\log Y}{\log T},
 \end{split}
\]
where the sums are enlarged to all primes. In the last inequality,
$U\ge\sqrt T$ and $\log(T/U)=\kappa\log Y$ were used.
Thus the remaining negative contribution in
\eqref{ss:average-identity} is
\[
 O\left(\frac{T(1+\kappa\log Y)}{L\log T}\right).
 \]
After division by $\vartheta(Q)T$, equation
\eqref{ss:prime-inputs} bounds this by
\[
 O_\kappa\left(\frac{(\log Y)^2}{L\log Z}\right)=o(1).
 \]
Here $\log Y\ge\log2$ absorbs the constant term in the numerator.
Subtracting this error from $43/64$ leaves at least $1/2$ for large
$M$, proving the lemma.
\end{proof}

\subsection{The general separated-support theorem}

\begin{theorem}[A uniform shared-hinge bound]
\label{ss:general}
Let $Y(M)\ge2$, $Z(M)\to\infty$, and integer functions $t(M)\ge0$,
$L(M)\ge1$ be prescribed. Assume
\begin{equation}
 \log Y=o(\sqrt{\log M}),\qquad
 \log Y=o(\log Z),\qquad
 \frac{(\log Y)^2}{L\log Z}\to0,\qquad
 C_*(t,L)\frac{\log Y}{\log Z}\to0.
 \label{ss:general-conditions}
\end{equation}
Let $Q\subseteq\{p:p\le Y\}$ and $R\subseteq\{p:p\ge Z\}$ be finite
prime sets. Let $A\subseteq[2,M]$ be distinct sources supported on
$Q\cup R$. For each source containing an $R$-prime, put
\[
 p(a)=P^+(a),\qquad d(a)=a/p(a).
\]
Let $\mathcal D$ be the distinct quotients and $R_0$ the union of their
$R$-prime supports. Suppose
\[
 \Omega_R(d)\le t\quad(d\in\mathcal D),
 \qquad \sum_{r\in R_0}\frac1r\le L/4.
\]
For every $\varepsilon>0$, all sufficiently large $M$ then satisfy,
uniformly over these data, all shifts, and all weights in $[0,1]$,
with $X(m)=\sum_{p\in Q\cup R}w_pv_p(m)$,
\begin{equation}
 \sum_{a\in A}(X(a)-2-\varepsilon)_+
 \le \sum_{b\in I}(X(b)-1)_+.
 \label{ss:shared-hinge}
\end{equation}
Consequently $\mathcal L_{\rm raw}(A,I)\le(2+\varepsilon)|A|$.
\end{theorem}

\begin{proof}
For large $M$, $Y<Z$, so $Q$ and $R$ are disjoint and the largest prime
of a mixed source belongs to $R$. Write $A_Q$ for the pure $Q$-sources
and $A_R=A\setminus A_Q$. If $A_R$ is empty, use
Proposition~\ref{ss:pure-core}.

Otherwise every $d=d(a)$ has $T_d\ge p(a)\ge Z$.
Every $R$-prime dividing $d$ is at most $p(a)\le T_d$, so (D) holds.
Let $m_d$ be the number of source preimages of $d$. Distinct preimages
have the form $dp(a)$ for distinct primes $p(a)\le T_d$, hence
$m_d\le\pi(T_d)$. Lemma~\ref{ss:average} and
\eqref{ss:prime-inputs} give
\begin{equation}
 \frac{m_d}{G_d}
 \le C_3\frac{\log Y}{\log Z},
 \label{ss:fibre-ratio}
\end{equation}
with an absolute constant, uniformly over all occupied fibres.

Since $w_{p(a)}\le1$,
$(X(a)-2-\varepsilon)_+\le(X(d)-1)_+=h_d$.
Using positivity of $G_d$, then Lemma~\ref{ss:capacity}, gives
\[
 \begin{split}
 \sum_{a\in A_R}(X(a)-2-\varepsilon)_+
 &\le\sum_d m_dh_d\\
 &\le C_3\frac{\log Y}{\log Z}\sum_dG_dh_d\\
 &\le C_3C_*(t,L)\frac{\log Y}{\log Z}
                  \sum_{b\in I}(X(b)-1)_+.
 \end{split}
\]
The last coefficient tends to zero.
For $V$ consisting only of the $Q$ weights,
Proposition~\ref{ss:pure-core} gives a uniform $\rho_\varepsilon<1$ with
\[
 \sum_{a\in A_Q}(X(a)-2-\varepsilon)_+
 \le\rho_\varepsilon H(V)\le\rho_\varepsilon H(X).
 \]
Increase the common threshold until the mixed coefficient is at most
$1-\rho_\varepsilon$, and add. This proves \eqref{ss:shared-hinge}.
Both source classes have used the same target hinge, rather than
independent primal covers.

Finally, the clipped raw boxed dual identity \cite{Synthesis} is
\[
 \mathcal L_{\rm raw}(A,I)=
 \max_{0\le w_p\le1}
 \left(\sum_{a\in A}X(a)-\sum_{b\in I}(X(b)-1)_+\right).
 \]
The weights are on the source primes. Using
$X(a)\le2+\varepsilon+(X(a)-2-\varepsilon)_+$ in this identity
proves the stated fractional bound.
\end{proof}

\subsection{Every fixed positive power cutoff}

\begin{corollary}
\label{ss:fixed-alpha}
Fix $\alpha>0$ and $Y(M)\ge2$ with $\log Y=o(\sqrt{\log M})$.
For distinct sources $A\subseteq[2,M]$ whose prime factors all satisfy
\[
 p\le Y(M)\quad\hbox{or}\quad p\ge M^\alpha,
\]
one has, uniformly over the source set and interval shift,
\[
 \mathcal L_{\rm raw}(A,I)\le(2+\varepsilon)|A|
 \quad\hbox{for all sufficiently large }M,
 \qquad
 \tau(A,I)\le(2+o(1))|A|\quad(|A|\to\infty).
\]
There is no separate bound on either the number of small-prime
occurrences or the number of large-prime occurrences in a source.
\end{corollary}

\begin{proof}
If $\alpha>1$, no source contains a large prime, and the pure-core
result suffices. Otherwise take $Z=M^\alpha$ and
$t=\lfloor1/\alpha\rfloor$. The size bound on each source ensures
$\Omega_R(d)\le t$.
If $r\in R_0$ occurs in $d(a)$, then $r\le p(a)$ and
$rp(a)\le a\le M$, so
$R_0\subseteq[M^\alpha,\sqrt M]$. Repeated prime occurrences are included
in this argument. For $0<\alpha<1/2$, Mertens' reciprocal estimate yields
\[
 \sum_{r\in R_0}\frac1r
 \le\log\frac1{2\alpha}+o(1).
 \]
For $\alpha>1/2$ the set is empty; at $\alpha=1/2$ at most the endpoint
prime occurs and its contribution is $o(1)$.
Choose a fixed integer $L$ with
\[
 L>4\left(1+\max\{0,\log(1/(2\alpha))\}\right).
 \]
The reciprocal condition then holds for large $M$.
All of \eqref{ss:general-conditions} follow from the assumed cutoff
on $Y$, since $t,L,C_*(t,L)$ are fixed. Theorem~\ref{ss:general}
proves the uniform raw bound.

The class is hereditary at the fixed ambient $M$: deleting sources
preserves their prime ranges. Lemma~\ref{lem:tr-hereditary}
therefore supplies the integer conclusion.
Its uniform large-$M$ threshold applies as $|A|\to\infty$, since
distinctness gives $M\ge|A|+1$.
\end{proof}

\subsection{A quantified cutoff below every fixed positive power}

The multiplicity estimate also allows $t$ to grow.
For a support restriction $p\ge Z$, take
$t=\lfloor\log M/\log Z\rfloor$ in the nonempty case $Z\le M$.
As before $R_0\subseteq[Z,\sqrt M]$.
When $Z\to\infty$, Mertens' reciprocal estimate gives
\[
 \sum_{r\in R_0}\frac1r\le\log t+o(1).
 \]
Indeed, if $Z\le\sqrt M$, the main term is
$\log((\log M)/(2\log Z))\le\log t$; otherwise $R_0$ is empty.
Thus the choice
\begin{equation}
 L=\lceil4\log(2t)+4\rceil
 \label{ss:growing-discount}
\end{equation}
satisfies the reciprocal condition for large $M$.

For $t\to\infty$, put $k=t+L-1$. Bounding the sum defining $C_2$ by
$L[k(t+1)]^{L-1}$ gives
\[
 \begin{split}
 \log C_*(t,L)
 &\le\log(t+L)+\log L
       +(L-1)\{\log(t+L-1)+\log(t+1)\}\\
 &\le8(\log t)^2+O(\log t).
 \end{split}
\]
The constant $8$ uses $L=4\log t+O(1)$ from
\eqref{ss:growing-discount}. The exponents of nonexceptional primes
are fixed by the target integer; assigning additional exponent choices
to them would lose this sharper estimate.

\begin{corollary}[An explicit growing-multiplicity class]
\label{ss:subpolynomial}
Fix $0<a<1/2$ and $c>0$ with $8c^2<1-a$. Put
\[
 Y(M)=\exp((\log M)^a),\qquad
 Z(M)=\exp\left(\frac{\log M}
                        {\exp(c\sqrt{\log\log M})}\right).
 \]
If every prime factor of every distinct source in $A\subseteq[2,M]$
is at most $Y(M)$ or at least $Z(M)$, then the uniform raw and
integer conclusions of Corollary~\ref{ss:fixed-alpha} hold.
\end{corollary}

\begin{proof}
Take $t=\lfloor\log M/\log Z\rfloor$ and $L$ from
\eqref{ss:growing-discount}. Then
$\log t=c\sqrt{\log\log M}+o(1)$, and the preceding bound gives
\[
 C_*(t,L)\le(\log M)^{8c^2+o(1)}.
 \]
Furthermore,
\[
 \frac{\log Y}{\log Z}=(\log M)^{a-1+o(1)},\qquad
 \frac{(\log Y)^2}{L\log Z}
       =(\log M)^{2a-1+o(1)}\longrightarrow0.
 \]
All conditions in \eqref{ss:general-conditions} hold, with
$C_*(t,L)\log Y/\log Z\to0$ following from $8c^2+a-1<0$.
Theorem~\ref{ss:general} proves the raw conclusion.
The support class is hereditary at fixed $M$, so
Lemma~\ref{lem:tr-hereditary} gives the integer conclusion.
\end{proof}

For the concrete choice $a=1/3$ and $c=1/4$,
\[
 C_*(t,L)\frac{\log Y}{\log Z}
 \le(\log M)^{-1/6+o(1)}\to0,\qquad
 \frac{(\log Y)^2}{L\log Z}
       =(\log M)^{-1/3+o(1)}\to0.
 \]
Here
$Z=M^{\exp(-\frac14\sqrt{\log\log M})}=M^{o(1)}$.
The theorem therefore does not require a fixed positive exponent
for the large-prime cutoff. It still imposes a substantial gap
between $Y$ and $Z$. The fixed-$\alpha$ theorem cannot be applied
with an arbitrary data-dependent $\alpha$ tending to zero; such
a choice requires the separate quantitative conditions above.
